\clearpage
\begin{landscape}

\chapter{Procedência das fontes}
\label{ap:procedencia}

\begin{promancha}
Todas as análises usam dados abertos e sem autenticação. As fontes não têm o
mesmo peso probatório, e a tabela final separa as revisadas por pares das
demais.
\end{promancha}

\section*{Reprodução}

\begin{promancha}
Cada figura exporta em CSV os dados que produzem cada painel, sob
\texttt{data/processed/}. Os notebooks e scripts geradores ficam em
\texttt{notebooks/}, e as figuras compostas em \texttt{figures/}. As métricas
citadas no corpo saem desses arquivos, e não da leitura da imagem.
\end{promancha}

\section*{Bases de dados}

\begin{promancha}
\campo{ONS, Portal de Dados Abertos} Restrição por \textit{constrained-off}
fotovoltaico e eólico; geração térmica por motivo de despacho; fator de
capacidade; curva de carga; intercâmbio nacional; energia armazenada por
subsistema; Custo Marginal de Operação semi-horário; relação entre conjunto e
usina; linhas de transmissão e subestações da Rede de Operação.

\campo{ANEEL} Sistema de Informações de Geração; relação de empreendimentos de
geração distribuída e informações técnicas fotovoltaicas; SIGET, com contrato,
módulo, linha de transmissão e subestações de origem e destino.

\campo{EPE} Webmap institucional, em serviço ArcGIS REST: camadas de linhas de
transmissão e subestações em operação, com geometria.

\campo{ESA e Copernicus} Sentinel-2 L2A, pelo catálogo STAC público da AWS em
\texttt{earth-\allowbreak search.\allowbreak aws.\allowbreak element84.\allowbreak com}.

\campo{NASA POWER} Irradiância diária, de MERRA-2 e SYN1DEG.

\campo{IBGE} Malhas territoriais das unidades da federação.
\end{promancha}

\section*{Comparação internacional}

\begin{promancha}
\tabfontw
\begin{tabular}{@{}>{\rag\arraybackslash}p{6.4cm}rr
                  >{\rag\arraybackslash}p{5.0cm}@{}}
\toprule
\cab{Sistema} & \cab{Taxa} & & \cab{Referência} \\
\midrule
Brasil, solar com apuração    & 21,6\% & & este relatório \\
Brasil, eólica e solar, 2025  & 20,7\% & & ONS \\
Chile, 2024                   & 17\%   & & imprensa setorial \\
Brasil, eólica com apuração   & 15,9\% & & este relatório \\
Irlanda, eólica, 2025         & 11,3\% & & EirGrid \\
Espanha, julho de 2025        & 11\%   & & imprensa setorial \\
Califórnia, solar             &  4,3\% & & Novan e Wang \\
\bottomrule
\end{tabular}
\notas{Os denominadores diferem, e a comparação com sistemas estrangeiros é de
ordem de grandeza, e não pareada. As duas linhas brasileiras são pareadas entre
si: ambas são razão de energia cortada sobre disponível, na janela de abril de
2024 a agosto de 2026, em que solar e eólica têm apuração simultânea, sobre 87
conjuntos solares e 178 eólicos. Numa extração única elas dão 22,1\% e 15,9\%; a
linha da solar preserva os 21,6\% da seção~\ref{sec:fig1}, de extração
anterior.}
\end{promancha}

\begin{promancha}
As taxas brasileiras ficam no topo da faixa internacional para as duas fontes. A
solar perde cerca de 1,4 vez o que a eólica perde. O que separa os dois regimes
não é a magnitude do corte, e sim a composição da causa: 91\% de origem
sistêmica na solar contra 76\% na eólica, medido na seção~\ref{sec:fig11}.
\end{promancha}

\section*{Mudança regulatória durante a execução}

\begin{promancha}
A pergunta que a seção~\ref{sec:fig3} deixou em aberto, sobre se o corte é
ressarcido, deixou de ser hipotética. A Lei 15.269/2025 criou compensação por
\textit{constrained-off} para eólicas e solares, custeada por Encargos de
Serviços do Sistema, com efeitos retroativos a 1 de setembro de 2023 e
condicionada à renúncia de ações judiciais e à assinatura de termo de
compromisso. A Portaria MME 140 regulamentou os cortes ocorridos até 25 de
novembro de 2025, e os posteriores ficam com a ANEEL.

A primeira fase reuniu manifestação de interesse de 1.539 usinas, com cerca de
53 GW instalados. Estimativas setoriais apontam perda da ordem de R\$ 6,5
bilhões em 2025, mesma ordem de grandeza da valoração ao CMO da
seção~\ref{sec:fig3} extrapolada para o parque.

Isso reforça a objeção daquela seção. O CMO mede valor sistêmico, e não
ressarcimento. Com a lei, os dois divergem por construção: o gerador é
compensado por energia cujo valor sistêmico, no instante do corte, era próximo
de zero.
\end{promancha}

\section*{Peso das fontes}

\begin{promancha}
\campo{Revisadas por pares} Vieira et al. (2026), Novan e Wang (2024),
Kruitwagen et al. (2021), Yu et al. (2018), Jiang et al. (2021), Alix-García e
Millimet (2023), Bird et al. (2016), Callaway e Sant'Anna (2021), Maji et al.
(2025).

\campo{Working papers e pré-prints} Newbery (2024), Proctor et al. (2023), Silva
et al. (2026), Global Renewables Watch (2025), Robinson et al. (2021).

\campo{Fontes normativas} Lei 15.269/2025, Portaria MME 140 e Lei 14.300/2022.

\campo{Imprensa setorial} Taxas de Chile, Espanha e Europa, e o valor de R\$ 6,5
bilhões. Usadas como ordem de grandeza, e nunca como evidência de mecanismo.
\end{promancha}

\end{landscape}
